\chapter{Caracteres y topología del grupo dual}\label{cap:v4-14}
Una frecuencia en la recta es la función $x\mapsto e^{i\xi x}$: transforma una suma de posiciones en un producto de números de módulo uno. Esa propiedad tiene sentido en cualquier grupo topológico y permite definir frecuencias sin coordenadas. El conjunto de esas frecuencias también tendrá una estructura de grupo y una topología. La lectura 25 de \cite{valencia4} proporciona el punto de partida; precisaremos la topología y las hipótesis de sus ejercicios.

Trabajamos con grupos de Hausdorff. Escribimos $\T=\{z\in\C:|z|=1\}$ y usamos notación multiplicativa, salvo en grupos como $\R$ y $\Z$. Un grupo \emph{LCA} es un grupo localmente compacto, abeliano y de Hausdorff. Para demostrar que el dual de un grupo LCA vuelve a ser localmente compacto necesitaremos el capítulo~\ref{cap:v4-15}. El ejercicio de exactitud requiere además la dualidad que se demostrará en el capítulo~\ref{cap:v4-16}; se indica esta dependencia expresamente para no esconder un teorema profundo bajo una verificación elemental.

\section{Frecuencias como homomorfismos}\label{lectura:v4-25}
\index{caracter@carácter!de un grupo}\index{grupo!dual}
\begin{definicion}
Un \emph{carácter} de $G$ es un homomorfismo continuo $\chi:G\to\T$. El conjunto de todos ellos se denota $\widehat G$. El producto de caracteres es puntual: $(\chi\eta)(x)=\chi(x)\eta(x)$.
\end{definicion}

\begin{proposicion}
El conjunto $\widehat G$ es un grupo abeliano. Su identidad es la función constante $1$ y $\chi^{-1}(x)=\overline{\chi(x)}=\chi(x^{-1})$. Sus elementos son exactamente las representaciones unitarias unidimensionales de $G$.
\end{proposicion}
\begin{proof}
El producto de funciones continuas es continuo, y
\[
 (\chi\eta)(xy)=\chi(x)\chi(y)\eta(x)\eta(y)
 =\chi(x)\eta(x)\chi(y)\eta(y).
\]
Se ha usado que los valores escalares conmutan, aunque $G$ no sea abeliano. La conjugación también preserva continuidad y la propiedad de homomorfismo. Las demás identidades del grupo se comprueban punto por punto en $\T$. Todo operador lineal sobre $\C$ es multiplicación por su valor en $1$, y preserva el producto interno si y sólo si ese valor tiene módulo uno. Por tanto una representación unitaria continua sobre $\C$ equivale precisamente a un carácter.
\end{proof}

\begin{lema}[No hay subgrupos pequeños en el círculo]\label{v4:arco-sin-subgrupos}
El único subgrupo de $\T$ contenido en $\{z:\operatorname{Re}z>0\}$ es $\{1\}$.
\end{lema}
\begin{proof}
Si $z=e^{i\theta}\ne1$, podemos elegir $0<|\theta|\le\pi$. Si $|\theta|\ge\pi/2$, el propio $z$ no pertenece al semiplano. Si $0<|\theta|<\pi/2$, elija el primer entero $n$ con $n|\theta|\ge\pi/2$. Entonces $\pi/2\le n|\theta|<\pi$, y $\operatorname{Re}(z^n)\le0$. Un subgrupo que contenga $z$ contiene esa potencia, de modo que no puede estar contenido en el semiplano indicado.
\end{proof}

\section{Convergencia uniforme sobre compactos}
\index{topologia@topología!compacto-abierta}
Para un espacio localmente compacto $X$, un compacto $K\subset X$ y un abierto $U\subset\C$, ponemos
\[
 [K,U]=\{f\in C(X):f(K)\subset U\}.
\]
La \emph{topología compacto-abierta} sobre $C(X)$ es la generada por estos conjuntos como subbase. La topología se define en el espacio de funciones $C(X)$, no en $X$; esto corrige una errata del enunciado original. La familia cubre $C(X)$, pues $[K,\C]=C(X)$ para cualquier $K$.

\begin{teorema}\label{v4:compacto-abierta}
El espacio $C(X)$ es Hausdorff con esta topología. Una red $f_\alpha$ converge a $f$ si y sólo si, para cada compacto $K\subset X$, $\sup_K|f_\alpha-f|\to0$.
\end{teorema}
\begin{proof}
Si $f\ne g$, existe $x$ con $f(x)\ne g(x)$. Dos discos abiertos disjuntos alrededor de esos valores producen vecindades disjuntas $[\{x\},U]$ y $[\{x\},V]$.

Supongamos convergencia uniforme en compactos y sea $f\in[K,U]$. Si $U=\C$ no hay nada que probar. En otro caso, el compacto $f(K)$ tiene distancia positiva al cerrado $\C\setminus U$: si esa distancia fuera cero, una sucesión en $f(K)$ con distancia tendiendo a cero tendría una subsucesión cuyo límite pertenecería simultáneamente a ambos conjuntos. Para un error uniforme menor que esa distancia, $f_\alpha(K)\subset U$. Esto vale para cada miembro de una intersección finita de subbásicos, y el conjunto de índices de una red permite elegir un umbral común para un número finito de condiciones.

Recíprocamente, sea $\eps>0$. La compacidad de $f(K)$ permite elegir $x_1,\ldots,x_m\in K$ cuyos discos de radio $\eps/4$ cubran $f(K)$. Los conjuntos
\[
 K_j=\{x\in K:|f(x)-f(x_j)|\le\eps/3\}
\]
son compactos, cubren $K$ y satisfacen $f(K_j)\subset B(f(x_j),2\eps/3)$. Así
\[
 \mathcal U=\bigcap_{j=1}^m[K_j,B(f(x_j),2\eps/3)]
\]
es una vecindad de $f$. Eventualmente $f_\alpha\in\mathcal U$. Para $x\in K_j$ obtenemos
$|f_\alpha(x)-f(x)|<2\eps/3+\eps/3=\eps$. Para asegurar una desigualdad estricta después de tomar supremos se repite el argumento con $\eps/2$. Resulta convergencia uniforme en $K$.
\end{proof}

En la lectura se intersectan ciertos abiertos con $K$ y se afirma que los conjuntos resultantes son compactos. Eso no es cierto en general. Los conjuntos cerrados $K_j$ de la prueba anterior evitan ese problema. Las redes permiten trabajar sin suponer una base numerable de vecindades: sus índices forman un conjunto dirigido, es decir, cada par tiene una cota superior, y «eventualmente» significa después de algún índice.

Dotamos a $\widehat G$ de la topología inducida por $C(G)$. Sus vecindades pueden describirse mediante errores uniformes en un número finito de compactos.
\begin{proposicion}\label{v4:dual-topologico}
El grupo $\widehat G$ es un grupo topológico abeliano y Hausdorff. Además, si $G,H$ son localmente compactos,
\[
 \widehat{G\times H}\simeq\widehat G\times\widehat H
\]
como grupos topológicos.
\end{proposicion}
\begin{proof}
Para caracteres, $|\chi(x)|=1$, y por tanto
\[
 |\chi(x)\overline{\eta(x)}-\chi_0(x)\overline{\eta_0(x)}|
 \le|\chi(x)-\chi_0(x)|+|\eta(x)-\eta_0(x)|.
\]
Tomar supremos en compactos demuestra continuidad conjunta del producto por el inverso; en particular da continuidad de producto e inversión. La separación Hausdorff se hereda de $C(G)$.

Para $\gamma\in\widehat{G\times H}$, las restricciones $\chi(x)=\gamma(x,e_H)$ y $\eta(y)=\gamma(e_G,y)$ son continuas y multiplicativas, y
$\gamma(x,y)=\chi(x)\eta(y)$. Recíprocamente, esa fórmula produce un carácter a partir de cualquier par $(\chi,\eta)$. Son aplicaciones inversas y homomorfismos. La restricción es continua porque $K\times\{e_H\}$ y $\{e_G\}\times L$ son compactos. La aplicación del par al producto es continua porque un compacto $S\subset G\times H$ tiene proyecciones compactas $K,L$, y el error uniforme del producto en $S$ está acotado por la suma de los errores de los factores en $K,L$.
\end{proof}

\begin{lema}[Compacidad de frecuencias y control común]\label{v4:dual-equicontinuidad}
Si $G$ es localmente compacto, la evaluación $(\chi,x)\mapsto\chi(x)$ en $\widehat G\times G$ es continua. Si $E\subset\widehat G$ es compacto, para todo $\eps>0$ existe una vecindad $U$ de $e$ tal que
\[
 \sup_{\chi\in E,\ x\in U}|\chi(x)-1|<\eps.
\]
\end{lema}
\begin{proof}
Fijados $(\chi_0,x_0)$, elija una vecindad de $x_0$ con clausura compacta $K$. En esa vecindad,
\[
 |\chi(x)-\chi_0(x_0)|\le\sup_K|\chi-\chi_0|+|\chi_0(x)-\chi_0(x_0)|.
\]
Se controla el primer término con una vecindad de $\chi_0$ y el segundo con una vecindad de $x_0$. Para la última afirmación, la continuidad en cada $(\chi,e)$ proporciona vecindades $V_\chi$ de $\chi$ y $U_\chi$ de $e$ donde el error es menor que $\eps/2$. Una subcubierta finita de los $V_\chi$ cubre $E$; la intersección de los correspondientes $U_\chi$ sirve simultáneamente para todo $E$.
\end{proof}

\section{Los modelos compacto, discreto y euclidiano}
\begin{teorema}\label{v4:dual-compacto-discreto}
Si $G$ es compacto, $\widehat G$ es discreto. Si $G$ es discreto, $\widehat G$ es compacto.
\end{teorema}
\begin{proof}
Si $G$ es compacto, el conjunto de caracteres cuya imagen está contenida en el semiplano derecho es una vecindad abierta de $1$ para la topología compacto-abierta. Esa imagen es un subgrupo de $\T$, luego el lema~\ref{v4:arco-sin-subgrupos} muestra que la vecindad contiene sólo el carácter trivial. Las traslaciones del grupo dual prueban que todos sus puntos son abiertos.

Si $G$ es discreto, toda función $G\to\T$ es continua y los compactos de $G$ son finitos. Por ello la topología compacto-abierta coincide con la convergencia puntual, esto es, con la topología inducida por el producto $\T^G$. El subconjunto de homomorfismos es cerrado: está definido por las ecuaciones cerradas $z_e=1$ y $z_{xy}=z_xz_y$, una por cada pareja. Tychonoff hace compacto a $\T^G$, y su cerrado $\widehat G$ es compacto. No se supone aquí que el dual sea discreto; en la lectura aparece esa palabra donde debe decir «$G$ es discreto».
\end{proof}

\begin{proposicion}\label{v4:duales-modelos}
Con sus topologías usuales,
\[
 \widehat\R\simeq\R,\quad \widehat\T\simeq\Z,\quad
 \widehat\Z\simeq\T,\quad
 \widehat{\Z/N\Z}\simeq\Z/N\Z.
\]
En el primer caso el parámetro $\xi$ corresponde a $x\mapsto e^{i\xi x}$.
\end{proposicion}
\begin{proof}
Para un carácter $\chi$ de $\R$, la continuidad en cero permite escribir $\chi(t)=e^{i\theta(t)}$ para $|t|<\delta$, con $\theta$ continua y $|\theta(t)|<\pi/3$. Siempre que $s,t,s+t$ pertenezcan a ese intervalo, el número $\theta(s+t)-\theta(s)-\theta(t)$ es un múltiplo de $2\pi$ y tiene valor absoluto menor que $\pi$, por lo que es cero. Esta aditividad local y continuidad implican $\theta(t)=\xi t$: fijando $0<a<\delta/2$, la subdivisión de $a$ en partes iguales da la fórmula en los múltiplos racionales de $a$ que permanecen en el intervalo, y la densidad de esos múltiplos da la fórmula en todo él. Para un $t$ cualquiera, elija $n$ con $t/n$ en ese intervalo; entonces $\chi(t)=\chi(t/n)^n=e^{i\xi t}$.

El parámetro es único, pues si $\xi\ne\eta$ basta elegir $t=\pi/(\xi-\eta)$ para distinguir los caracteres. La continuidad del parámetro hacia el dual se sigue de
$|e^{i\xi t}-e^{i\eta t}|\le|\xi-\eta||t|$. Para la continuidad inversa en cero, si $|\xi|\ge\eps$, el punto $t=\pi/|\xi|$ pertenece a $[-\pi/\eps,\pi/\eps]$ y $|e^{i\xi t}-1|=2$. Así un error uniforme menor que $1$ en ese compacto obliga a $|\xi|<\eps$. Las traslaciones de los dos grupos concluyen el argumento.

Un carácter de $\T$ compuesto con $t\mapsto e^{it}$ es $e^{i\xi t}$ y debe tener período $2\pi$. Por ello $\xi\in\Z$, y los caracteres son $z\mapsto z^n$. Su topología es discreta por compacidad de $\T$. La aplicación $t+2\pi\Z\mapsto e^{it}$ identifica $\R/(2\pi\Z)$ con $\T$: es una biyección continua, su dominio es compacto porque es imagen de $[0,2\pi]$, y el codominio es Hausdorff; por ello es un homeomorfismo. Esto recupera también el modelo del cociente empleado en la lectura. Un carácter de $\Z$ queda determinado por $z=\chi(1)$ y vale $\chi(n)=z^n$. La evaluación en $1$ es continua; la aplicación inversa es continua uniformemente en los subconjuntos finitos de $\Z$, luego es un homeomorfismo. Finalmente, los caracteres de $\Z/N\Z$ corresponden a los $z$ con $z^N=1$, que son $e^{2\pi i k/N}$, $0\le k<N$; sus productos suman los índices módulo $N$.
\end{proof}

\section{Ejercicios sobre productos, cocientes y duales}
\begin{ejercicio}[Lectura 25, ejercicio 1]\label{ej:v4-l25-1}
Sea $G$ un grupo abeliano finito de orden $n$. Demuestre que $\widehat G\simeq G$. Explique por qué la afirmación no proporciona, en general, una identificación canónica de $G$ con su dual.
\end{ejercicio}


\begin{ejercicio}[Lectura 25, ejercicio 2]\label{ej:v4-l25-2}
Sea $A\xrightarrow{f}B\xrightarrow{g}C$ una sucesión exacta corta de grupos LCA: $f$ identifica $A$ con $\ker g$ con su topología de subespacio, y $g$ identifica $C$ con $B/f(A)$ con su topología cociente. Pruebe que
\[
 \widehat C\xrightarrow{\widehat g}\widehat B
 \xrightarrow{\widehat f}\widehat A,
 \qquad \widehat g(\eta)=\eta\circ g,\quad
 \widehat f(\chi)=\chi\circ f,
\]
es exacta corta en el mismo sentido topológico.
\end{ejercicio}


\begin{ejercicio}[Lectura 25, ejercicio 3: topologías e hipótesis necesarias]\label{ej:v4-l25-3}
Demuestre las siguientes versiones precisas de los isomorfismos de la lectura:
\begin{enumerate}
\item Si cada $G_n$ es compacto, el dual de $\prod_{n\ge1}G_n$, con topología producto, es isomorfo a $\bigoplus_{n\ge1}\widehat G_n$, dotado de la topología discreta.
\item Si cada $G_n$ es discreto, el dual de $\bigoplus_{n\ge1}G_n$, dotado de la topología discreta, es isomorfo a $\prod_{n\ge1}\widehat G_n$, dotado de la topología producto.
\end{enumerate}
Explique por qué no basta suponer que los $G_n$ sean localmente compactos sin especificar las topologías de las sumas.
\end{ejercicio}


El grupo dual ya tiene operaciones continuas y una descripción concreta en los ejemplos básicos. El capítulo siguiente probará su compacidad local mediante el espectro del álgebra de convolución. Más adelante, la inversión y Plancherel permitirán recuperar el grupo original a partir de sus frecuencias y completar el ejercicio de exactitud.
